\documentclass[12pt,a4paper]{article}
\usepackage[utf8]{inputenc}
\usepackage[french]{babel}
\usepackage[T1]{fontenc}
\usepackage{geometry}
\geometry{left=2.5cm,right=2.5cm,top=2.5cm,bottom=2.5cm}
\usepackage{graphicx}
\usepackage{listings}
\usepackage{xcolor}
\usepackage{hyperref}
\usepackage{enumitem}
\usepackage{caption}
\usepackage{fancyhdr}

% Configuration des couleurs pour le code
\definecolor{codegreen}{rgb}{0,0.6,0}
\definecolor{codegray}{rgb}{0.5,0.5,0.5}
\definecolor{codepurple}{rgb}{0.58,0,0.82}
\definecolor{backcolour}{rgb}{0.95,0.95,0.92}

% Style pour le code C#
\lstdefinestyle{csharpstyle}{
    backgroundcolor=\color{backcolour},   
    commentstyle=\color{codegreen},
    keywordstyle=\color{magenta},
    numberstyle=\tiny\color{codegray},
    stringstyle=\color{codepurple},
    basicstyle=\ttfamily\footnotesize,
    breakatwhitespace=false,         
    breaklines=true,                 
    captionpos=b,                    
    keepspaces=true,                 
    numbers=left,                    
    numbersep=5pt,                  
    showspaces=false,                
    showstringspaces=false,
    showtabs=false,                  
    tabsize=2,
    language=[Sharp]C
}

\lstset{style=csharpstyle}

% En-tête et pied de page
\pagestyle{fancy}
\fancyhf{}
\rhead{Projet AR Tower Defense}
\lhead{Unity - AR Foundation}
\cfoot{\thepage}

\title{\textbf{Rapport de Projet} \\ \Large Application AR Tower Defense \\ avec Unity et AR Foundation}
\author{Projet de Réalité Augmentée}
\date{Mars 2026}

\begin{document}

\maketitle
\thispagestyle{empty}
\newpage

\tableofcontents
\newpage

\section{Introduction}

Ce projet consiste en une application de réalité augmentée (AR) développée avec Unity et AR Foundation. L'objectif est de créer un jeu de type \textit{tower defense} où des tourelles défendent une position contre des ennemis qui apparaissent dans l'environnement réel capturé par la caméra du smartphone.

\paragraph{Technologies utilisées :}
\begin{itemize}[label=$\bullet$]
    \item Unity 6000.3 (ou version similaire)
    \item AR Foundation
    \item ARCore XR Plugin (Android)
    \item C\# pour les scripts de gameplay
\end{itemize}

\section{Création du Projet - Étapes de Développement}

\subsection{Configuration Initiale}

\subsubsection{Installation des packages nécessaires}

Les packages suivants ont été installés via le Package Manager de Unity :

\begin{itemize}
    \item \textbf{AR Foundation} : Framework principal pour la réalité augmentée, fournissant une API unifiée
    \item \textbf{ARCore XR Plugin} : Support spécifique pour les appareils Android compatibles ARCore
    \item \textbf{Input System} : Système moderne de gestion des entrées Unity (remplace l'ancien système)
\end{itemize}

\subsubsection{Configuration du projet Android}

Dans \texttt{Project Settings > Player > Android} :

\begin{enumerate}
    \item \textbf{Minimum API Level} : Android 7.0 (API 24)
    \item \textbf{Target API Level} : Android 13 (API 33) ou supérieur
    \item \textbf{Graphics API} : Désactivation de "Auto Graphics API"
    \item Ajout manuel de Vulkan et OpenGLES3
    \item \textbf{Build System} : Gradle
    \item \textbf{Target Architectures} : ARM64 activé
\end{enumerate}

\subsection{Mise en Place de l'Environnement AR}

\subsubsection{Configuration de la scène principale}

\begin{enumerate}
    \item Suppression de la caméra par défaut de Unity
    \item Ajout d'un GameObject \texttt{AR Session} :
    \begin{itemize}
        \item Gère le cycle de vie de la session AR
        \item Initialise les services ARCore
    \end{itemize}
    \item Ajout d'un GameObject \texttt{AR Session Origin} contenant :
    \begin{itemize}
        \item \texttt{AR Camera} : Remplace la caméra classique, suit les mouvements du téléphone
        \item \texttt{AR Plane Manager} : Détecte les surfaces planes (sol, murs, tables)
        \item \texttt{AR Raycast Manager} : Permet de détecter où l'utilisateur touche dans l'espace AR
    \end{itemize}
\end{enumerate}

\subsubsection{Composants AR essentiels}

\paragraph{ARPlaneManager}
Configure pour détecter :
\begin{itemize}
    \item Plans horizontaux (sol, tables)
    \item Plans verticaux (murs) - optionnel
    \item Génère des maillages visuels pour les plans détectés
\end{itemize}

\paragraph{ARRaycastManager}
Essentiel pour le placement d'objets :
\begin{itemize}
    \item Convertit les touches écran en rayons 3D
    \item Détecte l'intersection avec les plans AR
    \item Retourne la position et rotation pour placement précis
\end{itemize}

\subsection{Création des Éléments de Jeu}

\subsubsection{Le système de placement}

Un script \texttt{PlacementController.cs} gère l'interaction utilisateur :

\begin{lstlisting}[caption=Extrait du PlacementController]
void Update() {
    if (Input.touchCount > 0) {
        Touch touch = Input.GetTouch(0);
        if (touch.phase == TouchPhase.Began) {
            List<ARRaycastHit> hits = new List<ARRaycastHit>();
            if (raycastManager.Raycast(touch.position, hits)) {
                Pose hitPose = hits[0].pose;
                Instantiate(prefab, hitPose.position, hitPose.rotation);
            }
        }
    }
}
\end{lstlisting}

\textbf{Fonctionnalités :}
\begin{itemize}
    \item Détection du touch sur l'écran
    \item Raycast vers les plans AR détectés
    \item Instanciation du prefab de tourelle à la position touchée
\end{itemize}

\subsubsection{La tourelle (Tower)}

\textbf{Composition :}
\begin{itemize}
    \item Modèle 3D simple (cylindre avec Material coloré)
    \item Script \texttt{TowerController.cs} avec :
    \begin{itemize}
        \item Système de détection d'ennemis dans un rayon défini
        \item Cadence de tir configurable
        \item LineRenderer pour l'effet visuel de tir
        \item Système de logs optionnel pour debug
    \end{itemize}
\end{itemize}

\textbf{Paramètres clés :}
\begin{lstlisting}[caption=Paramètres de la tourelle]
public float range = 0.05f;      // 5cm en metres reels
public float damage = 1f;         // Degats par tir
public float fireRate = 1f;       // 1 tir par seconde
public bool showLogs = false;     // Debug optionnel
\end{lstlisting}

\subsubsection{Les ennemis (Enemies)}

\textbf{Structure :}
\begin{itemize}
    \item Prefab avec tag "Enemy" (crucial pour la détection)
    \item Script \texttt{EnemyController.cs} gérant :
    \begin{itemize}
        \item Points de vie (health)
        \item Méthode \texttt{TakeDamage(float amount)}
        \item Destruction automatique quand la vie atteint 0
        \item Barre de vie visuelle (Canvas en World Space)
    \end{itemize}
    \item Spawn aléatoire autour d'une position définie
\end{itemize}

\subsubsection{Le gestionnaire de jeu (GameManager)}

Le \texttt{GameManager.cs} centralise :
\begin{itemize}
    \item Compteur d'ennemis tués
    \item Logs de progression du jeu
    \item Logique globale et état du jeu
    \item Pattern Singleton pour accès global
\end{itemize}

\section{Techniques et Choix de Conception}

\subsection{Système de Coordonnées AR}

\subsubsection{Problématique}
Les distances en AR sont exprimées en \textbf{mètres réels}, différentes des unités Unity classiques où 1 unité peut représenter n'importe quelle distance selon l'échelle du jeu.

\subsubsection{Solution adoptée}
\begin{itemize}
    \item Utilisation de petites valeurs pour \texttt{range = 0.05f} (5 centimètres)
    \item Ajustement de \texttt{lineWidth = 0.005f} pour la visibilité (5 millimètres)
    \item Tests itératifs pour trouver les bonnes proportions
    \item Documentation explicite des unités dans les commentaires
\end{itemize}

\subsection{Détection et Ciblage des Ennemis}

\subsubsection{Méthode choisie}
Utilisation de \texttt{GameObject.FindGameObjectsWithTag("Enemy")}

\subsubsection{Justification}
\begin{itemize}[label=$\checkmark$]
    \item Simple à implémenter et à comprendre
    \item Performant pour un petit nombre d'ennemis (< 50)
    \item Fiable avec le système de tags Unity natif
    \item Pas de setup complexe de colliders
\end{itemize}

\subsubsection{Alternative considérée}
\textit{Système de colliders sphériques avec triggers}
\begin{itemize}[label=$\times$]
    \item Plus performant pour beaucoup d'entités
    \item Requiert configuration de layers et collision matrix
    \item Complexité inutile pour ce prototype
\end{itemize}

\subsection{Feedback Visuel}

\subsubsection{LineRenderer pour les tirs}

Configuration du LineRenderer :
\begin{lstlisting}[caption=Configuration visuelle du tir]
shootLine.startColor = Color.yellow;
shootLine.endColor = Color.red;
shootLine.startWidth = 0.005f;  // 5mm
shootLine.endWidth = 0.005f;
shootLine.material.shader = Shader.Find("Sprites/Default");
\end{lstlisting}

\textbf{Avantages :}
\begin{itemize}
    \item Feedback immédiat de l'action de tir
    \item Peu coûteux en performance (pas de particules)
    \item Effet visuel clair et lisible en AR
    \item Système de fondu progressif pour effet naturel
\end{itemize}

\subsubsection{Système de fade}

\begin{lstlisting}[caption=Coroutine de fondu]
IEnumerator FadeOutLine() {
    yield return new WaitForSeconds(0.1f);
    shootLine.enabled = false;
}
\end{lstlisting}

Permet une disparition progressive en 0.1 seconde pour un effet fluide sans surcharge.

\section{Problèmes Rencontrés et Solutions}

\subsection{Portée de Tir Trop Grande}

\subsubsection{Symptôme}
La tourelle tirait sur des ennemis visuellement très éloignés, rendant le jeu sans difficulté.

\subsubsection{Cause}
Valeur initiale \texttt{range = 5f} adaptée aux jeux 3D classiques mais représentant 5 mètres en réalité augmentée (distance excessive).

\subsubsection{Solution}
\begin{enumerate}
    \item Réduction drastique : \texttt{range = 0.05f} (5 centimètres)
    \item Ajout de logs détaillés pour afficher les distances réelles :
\begin{lstlisting}
Debug.Log($"[Tower] Ennemi a distance: {distance:F3}, Range: {range:F6}");
\end{lstlisting}
    \item Tests progressifs avec différentes valeurs : 0.1f, 0.05f, 0.03f
    \item Documentation du ratio distance/gameplay
\end{enumerate}

\subsection{Script Manquant (Warning)}

\subsubsection{Message d'erreur}
\texttt{The referenced script (Unknown) on this Behaviour is missing!}

\subsubsection{Causes probables}
\begin{itemize}
    \item Script supprimé mais encore référencé sur un GameObject
    \item Component UI avec script manquant après refactoring
    \item Prefab avec référence cassée
\end{itemize}

\subsubsection{Impact}
Minimal sur le jeu principal, avertissement seulement.

\subsubsection{Résolution}
\begin{enumerate}
    \item Recherche des GameObjects avec composants manquants
    \item Nettoyage des références cassées dans les prefabs
    \item Validation de tous les scripts dans la scène
\end{enumerate}

\subsection{Visibilité du LineRenderer}

\subsubsection{Problème initial}
Ligne de tir invisible ou trop fine, ne s'affichant pas correctement en AR.

\subsubsection{Ajustements effectués}
\begin{lstlisting}
shootLine.startWidth = 0.005f;  // 5mm en AR (visible mais pas envahissant)
shootLine.endWidth = 0.005f;
shootLine.material.shader = Shader.Find("Sprites/Default");
\end{lstlisting}

\subsubsection{Shader choisi}
\texttt{Sprites/Default} pour compatibilité garantie sur mobile et rendu correct en AR.

\subsection{Performance sur Mobile}

\subsubsection{Problématiques identifiées}
\begin{itemize}
    \item Boucle de détection appelée chaque frame
    \item LineRenderer actif en permanence
    \item Logs verbeux réduisant les FPS
\end{itemize}

\subsubsection{Optimisations appliquées}
\begin{enumerate}
    \item Limitation de la fréquence de tir (\texttt{fireRate = 1f})
    \item Break après ciblage du premier ennemi (pas de multi-target)
    \item Désactivation du LineRenderer quand non utilisé
    \item Logs conditionnels avec flag \texttt{showLogs = false}
    \item Pas de calcul de distance superflue
\end{enumerate}

\subsection{Détection des Plans AR Instable}

\subsubsection{Symptômes}
\begin{itemize}
    \item Placement d'objets difficile ou impossible
    \item Plans non détectés malgré surfaces planes
    \item Grille AR clignotante
\end{itemize}

\subsubsection{Solutions environnementales}
\begin{enumerate}
    \item \textbf{Éclairage} : Bonne luminosité ambiante (pas de contre-jour)
    \item \textbf{Textures} : Environnement avec motifs distincts
    \begin{itemize}
        \item Éviter surfaces blanches uniformes
        \item Privilégier sols avec motifs
    \end{itemize}
    \item \textbf{Mouvement} : Déplacer lentement le téléphone pour scanner
    \item \textbf{Stabilité} : Garder l'appareil stable après détection
\end{enumerate}

\subsubsection{Solutions techniques}
Configuration \texttt{ARPlaneManager} :
\begin{itemize}
    \item Activer détection horizontale ET verticale
    \item Ajuster \texttt{Detection Mode} selon besoin
    \item Valider les permissions caméra dans AndroidManifest
\end{itemize}

\section{Utilisation de l'Application Finale}

\subsection{Installation}

\subsubsection{Prérequis}
\begin{itemize}
    \item Appareil Android avec ARCore (Android 7.0+)
    \item Connexion USB activée (mode développeur)
    \item Drivers ADB installés sur PC
\end{itemize}

\subsubsection{Procédure d'installation}
\begin{enumerate}
    \item Connecter l'appareil Android au PC
    \item Dans Unity : \texttt{File > Build Settings}
    \item Sélectionner Android, cliquer \texttt{Build and Run}
    \item Accepter les permissions caméra au premier lancement
\end{enumerate}

\textit{Alternative} : Installer l'APK généré via USB ou ARCore.

\subsection{Lancement du Jeu}

\subsubsection{Étape 1 : Scanner l'environnement}
\begin{enumerate}
    \item Orienter le téléphone vers le sol ou une table
    \item Bouger lentement de gauche à droite
    \item Attendre l'apparition de la grille AR blanche
    \item Continuer jusqu'à avoir une zone stable
\end{enumerate}

\subsubsection{Étape 2 : Placer la tourelle}
\begin{enumerate}
    \item Toucher l'écran sur un plan détecté
    \item La tourelle apparaît instantanément
    \item \textit{Note} : Une seule tourelle pour cette version
\end{enumerate}

\subsubsection{Étape 3 : Observer le gameplay}
\begin{itemize}
    \item Les ennemis spawn automatiquement autour de la tourelle
    \item La tourelle tire automatiquement dans sa portée
    \item Effet visuel : ligne jaune-rouge lors du tir
    \item Observer le compteur dans les logs (si activé)
\end{itemize}

\subsection{Conseils d'Utilisation}

\subsubsection{Pour une meilleure expérience}
\begin{itemize}[label=$\star$]
    \item \textbf{Éclairage} : Utiliser dans un environnement bien éclairé
    \item \textbf{Surfaces} : Éviter surfaces réfléchissantes, transparentes ou noires
    \item \textbf{Stabilité} : Garder l'appareil relativement stable après placement
    \item \textbf{Redémarrage} : Fermer et relancer l'app si la détection échoue
    \item \textbf{Batterie} : L'AR consomme beaucoup, garder le téléphone chargé
\end{itemize}

\section{Architecture du Code}

\subsection{Séparation des Responsabilités}

\subsubsection{Principe appliqué}
Chaque script a une responsabilité unique et bien définie (Single Responsibility Principle).

\begin{table}[h]
\centering
\begin{tabular}{|l|p{8cm}|}
\hline
\textbf{Script} & \textbf{Responsabilité} \\
\hline
PlacementController & Interface utilisateur (touch), Interaction AR, Instanciation \\
\hline
TowerController & Logique de combat, Détection ennemis, Effets visuels \\
\hline
EnemyController & Gestion de vie, Réception dégâts, Comportement ennemi \\
\hline
GameManager & État global, Statistiques, Coordination \\
\hline
\end{tabular}
\caption{Répartition des responsabilités}
\end{table}

\subsection{Paramètres Exposés}

Tous les scripts exposent leurs paramètres clés dans l'Inspector Unity :

\begin{lstlisting}[caption=Exemple de paramètres publics]
public class TowerController : MonoBehaviour {
    public float range = 0.05f;      // Modifiable sans recompiler
    public float fireRate = 1f;       // Ajustable en temps reel
    public bool showLogs = false;     // Debug on/off
}
\end{lstlisting}

\subsubsection{Avantages}
\begin{itemize}
    \item Ajustement des valeurs sans recompilation
    \item Tests rapides de différentes configurations
    \item Accessibilité pour les game designers non-programmeurs
    \item Itération rapide du gameplay
\end{itemize}

\subsection{Gestion des Events}

\subsubsection{Pattern Observer}
Utilisation de méthodes callback pour les événements :
\begin{lstlisting}
// Quand un ennemi meurt
GameManager.Instance.OnEnemyKilled();
\end{lstlisting}

\subsubsection{Avantages}
\begin{itemize}
    \item Couplage faible entre composants
    \item Facilite l'ajout de nouveaux comportements
    \item Maintenabilité améliorée
\end{itemize}

\section{Améliorations Futures}

\subsection{Gameplay}

\subsubsection{Ennemis mobiles}
\begin{itemize}
    \item Implémentation d'un système de pathfinding (A*)
    \item Déplacement vers un objectif à défendre
    \item Détection d'obstacles AR
    \item Animation de marche/course
\end{itemize}

\subsubsection{Système de vagues}
\begin{itemize}
    \item Spawn progressif par vagues
    \item Difficulté croissante (vitesse, nombre, résistance)
    \item Pause entre les vagues
    \item Boss en fin de vague
\end{itemize}

\subsubsection{Types d'ennemis variés}
\begin{itemize}
    \item Ennemis rapides / lents
    \item Ennemis résistants / fragiles
    \item Ennemis volants
    \item Ennemis avec capacités spéciales
\end{itemize}

\subsubsection{Tourelles multiples}
\begin{itemize}
    \item Placement illimité avec système de coût
    \item Ressources gagnées en tuant des ennemis
    \item Différents types de tourelles (rapide, lente puissante, zone)
    \item Amélioration/upgrade des tourelles
\end{itemize}

\subsection{Interface Utilisateur}

\subsubsection{UI Canvas en AR}
\begin{itemize}
    \item Affichage des stats en temps réel (argent, vague, santé)
    \item Interface en World Space attachée à l'AR
    \item Mini-map des ennemis et tourelles
    \item Indicateurs de portée des tourelles
\end{itemize}

\subsubsection{Boutons de contrôle}
\begin{itemize}
    \item Pause / Resume du jeu
    \item Restart après défaite
    \item Menu de sélection de tourelle
    \item Zoom et navigation caméra
\end{itemize}

\subsubsection{Feedback audio}
\begin{itemize}
    \item Sons de tir différenciés par type
    \item Explosion à la destruction d'ennemis
    \item Musique d'ambiance
    \item Alertes sonores (vague, défaite)
\end{itemize}

\subsection{Technique}

\subsubsection{Object Pooling}
\begin{itemize}
    \item Réutilisation des GameObjects ennemis
    \item Pool de projectiles
    \item Réduction du Garbage Collection
    \item Amélioration performance significative
\end{itemize}

\subsubsection{Occlusion AR}
\begin{itemize}
    \item Ennemis cachés derrière objets réels
    \item Utilisation de Depth API d'ARCore
    \item Réalisme augmenté de l'intégration AR
\end{itemize}

\subsubsection{Persistance et Cloud Anchors}
\begin{itemize}
    \item Sauvegarde de position des tourelles
    \item AR Cloud Anchors pour multi-session
    \item Partage du terrain entre joueurs
    \item Mode multijoueur coopératif
\end{itemize}

\subsubsection{Analytics et métriques}
\begin{itemize}
    \item Tracking des sessions de jeu
    \item Statistiques de performance
    \item Heatmap de placement de tourelles
    \item Optimisation based on data
\end{itemize}

\section{Conclusion}

\subsection{Bilan du Projet}

Ce projet a permis d'explorer les bases de la réalité augmentée avec Unity AR Foundation. Les principaux défis rencontrés ont été :

\begin{enumerate}
    \item \textbf{Compréhension de l'échelle AR} : Adaptation des distances et proportions du jeu classique vers les mètres réels de l'AR
    \item \textbf{Configuration mobile} : Gestion des spécificités Android, ARCore, et permissions
    \item \textbf{Debug sur appareil} : Difficulté à débugger sans connection directe, utilisation intensive de logs et tests itératifs
    \item \textbf{Performance mobile} : Optimisation constante pour maintenir 30 FPS minimum
\end{enumerate}

\subsection{Compétences Acquises}

\subsubsection{Techniques AR}
\begin{itemize}[label=$\checkmark$]
    \item Setup complet d'un projet AR Foundation
    \item Détection de plans et raycast dans l'espace AR
    \item Placement d'objets virtuels dans le monde réel
    \item Gestion du tracking et de la stabilité AR
\end{itemize}

\subsubsection{Développement Unity}
\begin{itemize}[label=$\checkmark$]
    \item Architecture de code modulaire et maintenable
    \item Utilisation avancée de LineRenderer
    \item Système de tags et détection d'objets
    \item Coroutines pour effets temporels
\end{itemize}

\subsubsection{Développement Mobile}
\begin{itemize}[label=$\checkmark$]
    \item Build et déploiement Android
    \item Optimisation pour appareils mobiles
    \item Debug et logs sur device
    \item Gestion des inputs tactiles
\end{itemize}

\subsection{Points Forts du Projet}

\begin{itemize}[label=$\star$]
    \item \textbf{Fonctionnel} : Système de base opérationnel (placement + tir automatique)
    \item \textbf{Visuel} : Feedback clair et compréhensible
    \item \textbf{Code} : Structure modulaire et paramétrable
    \item \textbf{Documentation} : Logs détaillés et commentaires
    \item \textbf{Extensible} : Base solide pour évolutions futures
\end{itemize}

\subsection{Perspective}

Le projet constitue une \textbf{base solide} pour un tower defense AR plus complet, avec toutes les fondations techniques en place pour l'expansion. Les points critiques (détection AR, placement, ciblage, tir) sont fonctionnels et optimisés.

Les prochaines étapes logiques seraient :
\begin{enumerate}
    \item Implémentation du pathfinding ennemi
    \item Système économique et placement multiple
    \item UI complète et polish visuel
    \item Tests utilisateurs et balancing gameplay
\end{enumerate}

L'expérience acquise lors de ce projet est directement transférable à d'autres applications AR : visualisation de produits, jeux AR avancés, applications éducatives, ou outils professionnels de réalité augmentée.

\vspace{1cm}

\begin{center}
\textit{Fin du rapport}
\end{center}

\end{document}
